\section{Discussion}
\label{sec:analysis}

\subsection{Cross-Domain Transfer}
\label{sec:cross_domain}

To understand whether visual tool-use behavior transfers across domains, we fine-tune Qwen3.5-9B on each single-domain slice of \dataset and evaluate each model on all five domains. Table~\ref{tab:cross} shows that single-domain training often improves out-of-domain performance, but the transfer is highly asymmetric. For example, GUI-only training raises Visual Search from 35.6 to 52.8, and Visual-Search-only training raises Table from 25.8 to 40.9. These gains suggest that some behaviors, such as inspecting localized image regions and grounding visual evidence before answering, are reusable across domains.

At the same time, no single domain is a substitute for the full corpus. \dataset obtains the best average score, and it is best on four of the five evaluation domains. GUI is the exception: the GUI-only slice reaches 36.3, while the full corpus reaches 31.6. This mirrors the no-tool trajectory result in Table~\ref{tab:notool} and indicates that GUI Grounding is unusually sensitive to the training mixture and action format.

\begin{table*}[t]
\centering
\small
\setlength{\tabcolsep}{5pt}
\begin{tabular}{lrrrrrr}
\toprule
Train slice (Qwen3.5-9B) & Chart & GUI & Table & Vis.\ Search & Web2HTML & AVG \\
\midrule
w/o tool          & 34.6 & 27.1 & 25.8 & 35.6 & 40.1 & 32.6 \\
\midrule
Chart             & 46.6 & 22.2 & 37.1 & 44.3 & 36.5 & 37.4 \\
GUI               & 41.4 & \textbf{36.3} & 38.7 & 52.8 & 39.1 & 41.7 \\
Table             & 45.7 & 20.5 & 38.9 & 45.8 & 39.2 & 38.0 \\
Visual Search     & 45.7 & 23.3 & 40.9 & 57.3 & 21.6 & 37.8 \\
Web2HTML          & 41.6 & 18.2 & 38.5 & 42.2 & 44.6 & 37.0 \\
\midrule
\dataset          & \textbf{49.3} & 31.6 & \textbf{43.1} & \textbf{59.4} & \textbf{45.5} & \textbf{45.8} \\
\bottomrule
\end{tabular}
\caption{Cross-domain transfer. Each single-domain row trains on one slice of \dataset for three epochs and evaluates on all domains. Bold marks the best score in each evaluation column. The full \dataset row is shown for reference.}
\label{tab:cross}
\end{table*}

\subsection{Interpreting the GUI Exception}
\label{sec:gui_exception}

GUI Grounding is the only domain where text-only trajectory SFT is stronger than OpenVisTool tool-use SFT: 40.0 versus 31.6. We report this directly because it clarifies what the current recipe does and does not solve. The result does not mean that tool-use supervision is useless for GUI: \dataset still improves the raw Qwen3.5-9B baseline from 27.1 to 31.6, and the GUI-only tool-use slice reaches 36.3. Rather, the result suggests that our current GUI tool protocol is less sample-efficient than direct point imitation under a point-in-box metric.

There are two likely causes. First, GUI Grounding has a very narrow output space: the model ultimately needs to emit one coordinate. A no-tool trajectory can train this behavior directly, whereas a tool-use trajectory introduces extra decisions such as where to crop, how to map crop-local evidence back to the original screenshot, and when to issue the final \texttt{computer\_use} action. Second, cross-domain training appears to dilute GUI-specific action patterns. Most non-GUI single-domain slices reduce GUI performance below the raw baseline, and the GUI-only slice outperforms the full mixture on GUI. This points to GUI-specific interface design and mixture balancing as important future improvements.

\subsection{What the Filtering Ablation Shows}
\label{sec:filtering_discussion}

The filtering ablation supports the central data-quality claim. Correctness alone produces a strong corpus, but it can retain trajectories where the final answer is right even if the tool calls are incidental. Tool-gain filtering addresses that failure mode, and adding it to the correctness filter improves every domain despite reducing the corpus to 42K examples. Conversely, tool-gain-only is weaker than correctness-only, showing that a useful-looking tool prefix is not enough if the teacher trajectory itself is unreliable.

The effect is largest on Table and Visual Search. Compared with correctness-only, the full \dataset gains +4.2 on Table and +2.1 on Visual Search; compared with tool-gain-only, it gains +5.6 and +4.7. These are the domains where the model most needs to acquire evidence through tool outputs, so they benefit most from retaining trajectories that are both correct and causally useful.

\subsection{Takeaway}
\label{sec:analysis_takeaway}

The results support a nuanced view of visual tool-use distillation. Tool-use trajectories are most valuable for evidence-acquisition tasks, where the tool output reveals information that the model cannot reliably infer from the original input. This explains the large gains on Visual Search, Table, and Chart, as well as the advantage over no-tool trajectory SFT on four of five domains. GUI Grounding exposes a limitation: when the benchmark rewards a single precise action, direct imitation can outperform the current multi-step tool-use protocol. We therefore frame OpenVisTool as a strong cross-domain recipe for visual tool-use data, while treating GUI as a diagnostic case for better action interfaces and domain-aware mixing.
